\section{Proofs for Section~\ref{sec:fields}}
\label{app:fields}

Throughout this appendix, $\norm{\cdot}$ is the Euclidean norm of a vector
and the operator norm of a matrix, and $\norm{\cdot}_F$ is the Frobenius norm.
For a quadratic $g(X)=g_0+b^{\mathsf T}X+X^{\mathsf T}TX$ with $T$
symmetric, $\Gamma(g)$ is its canonical Hessian Gram
\eqref{eq:algebraic-canonical}, with blocks $2bb^{\mathsf T}+4g_0T$,
$\Delta(b,T)$, and $\Xi(T)$, and for a center $c\in\R^n$ we write
$\Gamma_c(g)=\Gamma(g(\cdot+c))$ for the canonical Gram of the recentred
quadratic; both are extended linearly to signed rational combinations
$\sum_j\sigma_jg_j^2$. By Lemma~\ref{lem:algebraic-gram},
$v^{\mathsf T}\nabla^2\bigl(\sum_j\sigma_jg_j^2\bigr)(c+u)v
=(v,u\otimes v)^{\mathsf T}\Gamma_c(\cdot)(v,u\otimes v)$, and
$\Gamma=\Psi_c^{\mathsf T}\Gamma_c\Psi_c$ with the invertible matrix $\Psi_c$
of that lemma. Two consequences are used repeatedly. If $\Gamma_c$ of a
combination is positive definite, then the matrix $\Gamma$ computed from the
same factors at the origin is a positive definite full Hessian Gram; it is
rational when the factors and weights are rational, and an integer matrix
when the factors are integer quadratics and the weights are integers, since
$g_0$, $b$, and $2T$ are then integral and every block of
\eqref{eq:algebraic-canonical} is an integer combination of their entries.
Second, $\norm{\Delta(b,T)}\le\norm{\Delta(b,T)}_F\le6\norm b\norm T_F$, the
two tensors in its definition having Frobenius norms $2\norm b\norm T_F$ and
$4\norm b\norm T_F$; $\norm{\Xi(T)}\le8\norm T_F^2+4\norm T^2$; and for a
symmetric block matrix
$\norm{\left(\begin{smallmatrix}C&D\\D^{\mathsf T}&Q\end{smallmatrix}\right)}
\le\max\{\norm C,\norm Q\}+\norm D$.

Lemma~\ref{lem:quartic-realization} is used with $m=n$ residuals, in the
notation of its statement: quadratics $g,r_1,\ldots,r_n$ vanishing at $p$,
positive rationals $\mu,\Lambda,\beta,\nu$ satisfying
\eqref{eq:algebraic-realization-data}, and a rational square $\varepsilon=t^2$
satisfying \eqref{eq:algebraic-realization-eps}. Its conclusions (b) and (c)
give, for $f=(g/(t\nu))^2+\sum_j(r_j/\nu)^2$, the bound
$\nabla^2f\succeq\frac32I$ and a rational canonical Gram
$\Gamma_f\succeq\gamma I$ with
$\gamma=\min\{3/2,2\mu^2/(\varepsilon\nu^2)\}/
\bigl(((1+\theta)^2+1)(1+\norm p)^2\bigr)$ and
$\theta=3\sqrt n\,\varepsilon(\Lambda+n\beta)/\mu^2$.

\subsection{Proof of Lemma~\ref{lem:fields-tools}}
\label{app:fields-tools}

(a) Let $\delta$ be the largest degree of a factor and suppose $\delta\ge3$.
The homogeneous part of degree $2\delta$ of $\sum_jg_j^2$ is
$\sum_j(g_j^{(\delta)})^2$, where $g_j^{(\delta)}$ is the degree-$\delta$
part of $g_j$. It is positive at any real point where some
$g_j^{(\delta)}$ is nonzero, so it is a nonzero polynomial of degree
$2\delta>4$, which is impossible.

(b) The real numbers $g_j(p)$ satisfy $\sum_jg_j(p)^2=f(p)=0$. For the Gram
statement, $z(p)^{\mathsf T}Qz(p)=0$ and $Q\succeq0$ give
$Q^{1/2}z(p)=0$, hence $Qz(p)=0$. The $i$th entry of $Qz(p)$ is the value
at $p$ of the quadratic whose coefficient vector is the $i$th column of $Q$.
If $Q$ has entries in $E$, every column of $Q$ is therefore the coefficient
vector of an element of $\mathcal I_2^E(p)$. Let $B$ be the matrix with
entries in $E$ whose columns are the coefficient vectors of the basis $g$,
so that $g=B^{\mathsf T}z$; it has full column rank, and we fix a left
inverse $C$ with entries in $E$. Every column of $Q$ lies in the range of
$B$, so $BCQ=Q$, and by symmetry $Q=QC^{\mathsf T}B^{\mathsf T}$. Thus
$Q=B(CQC^{\mathsf T})B^{\mathsf T}$, and $S=CQC^{\mathsf T}$ is positive
semidefinite with entries in $E$. Finally
$f=z^{\mathsf T}Qz=g^{\mathsf T}Sg$.

(c) If $f=\sum_jg_j^2$ with $g_j\in E[X]$, parts (a) and (b) put every
$g_j$ in $\mathcal I_2^E(p)$, so $\varphi(f)=\sum_j\varphi(g_j^2)\ge0$. If
$f$ has a positive semidefinite Gram over $E$, part (b) gives
$f=g^{\mathsf T}Sg$ with $S\succeq0$. Let $\Phi_{ij}=\varphi(g_ig_j)$.
For $c\in E^t$ we have $c^{\mathsf T}\Phi c=\varphi((c^{\mathsf T}g)^2)\ge0$.
Since $\Q\subseteq E$ is dense in $\R$ and $\Phi$ is a real matrix,
$\Phi\succeq0$. Therefore $\varphi(f)=\tr(S\Phi)\ge0$. Both cases contradict
$\varphi(f)<0$.

(d) Let $P_a$ be the minimal polynomial of $a$ over $E$ and $r=\deg P_a$.
It divides $T^D-2$, whose roots are $a\zeta^j$ with $\zeta=e^{2\pi\mathrm
i/D}$. The constant term of $P_a$ is $(-1)^ra^r\zeta^e$ for some integer
$e$, and it is real. The only real $D$th roots of unity are $\pm1$, and
$-1$ is not one because $D$ is odd, so $\zeta^e=1$ and $a^r\in E$. Let
$d'=\gcd(r,D)=ur+vD$ with integers $u,v$. Then
$a^{d'}=(a^r)^u(a^D)^v=(a^r)^u2^v\in E$, so $a$ is a root of
$T^{d'}-a^{d'}\in E[T]$, and minimality gives $r\le d'$. Hence $r=d'$
divides $D$, and $T^r-a^r$ is the minimal polynomial. If $D=\ell$ is prime and
$a\notin E$, then $r>1$ forces $r=\ell$. If $D=5^k$ and $a\notin E$, then
$r=5^j$ with $j\ge1$. The element $a^5$ is a root of
$T^{r/5}-a^r\in E[T]$, so $[E(a^5):E]\le r/5$, while
$[E(a):E(a^5)]\le5$ because $a$ is a root of $T^5-a^5$. Since
$[E(a):E]=r$, both inequalities are equalities.

\subsection{The prime family}
\label{app:fields-prime}

Fix a prime $\ell\ge5$, $s=(\ell-1)/2\ge2$, $n=s+1\ge3$, $a=2^{1/\ell}$, and
$p=(a^s,\ldots,a^{2s})$. All coordinates of $p$ lie in $[1,2)$, and
$a<2^{1/5}<1.15$.

\paragraph{An exposing quadratic.}
Define the integer quadratics
\[
 P_0=x_0^2-x_s,\quad P_1=x_s^2-2x_{s-1},\quad P_2=x_0x_1-2,\quad
 P_{i+2}=x_i^2-x_{i-1}x_{i+1}\quad(1\le i\le s-1),
\]
the real weights $w_0=2a$, $w_1=a^2$, $w_2=-2$,
$w_{i+2}=a^{2s-2i+2}$ ($1\le i\le s-1$), and
$E_*=\sum_{j=0}^{s+1}w_jP_j$. Every $P_j$ vanishes at $p$, because
$a^{2s+1}=2$. None of them has a linear term in $x_0$; for $P_1$ this uses
$s\ge2$. Let $\Theta$ be the symmetric tridiagonal matrix indexed by
$0,\ldots,s$ with $\Theta_{ii}=a^{2s-2i}$ and
$\Theta_{i,i+1}=\Theta_{i+1,i}=\frac12a^{2s-2i-1}$, and put
\[
 L(x)=(x_1-ax_0,\ x_2-ax_1,\ \ldots,\ x_s-ax_{s-1},\ 2-ax_s)^{\mathsf T}.
\]
We claim $E_*=L^{\mathsf T}\Theta L$. Write $L_i$ for the entries. Using
$a^{2s+1}=2$, the coefficients of $\sum_ia^{2s-2i}L_i^2+
\sum_{i<s}a^{2s-2i-1}L_iL_{i+1}$ are as follows. The coefficient of $x_0^2$
is $a^{2s+2}=2a$. For $1\le j\le s-1$ the coefficient of $x_j^2$ is
$a^{2s-2j+2}(1+1-1)=a^{2s-2j+2}$, from $L_{j-1}^2$, $L_j^2$, and
$L_{j-1}L_j$; likewise the coefficient of $x_s^2$ is $a^2(1+1-1)=a^2$. The
coefficient of $x_jx_{j+1}$ is $-2a^{2s-2j+1}+a^{2s-2j+1}+a^{2s-2j+1}=0$ for
$1\le j\le s-1$, and $-2a^{2s+1}+a^{2s+1}=-2$ for $j=0$. The coefficient of
$x_jx_{j+2}$ is $-a^{2s-2j}$, from $L_jL_{j+1}$. The linear terms are $-2ax_s$
and $-2a^2x_{s-1}$, and the constant is $4$. These are exactly the
coefficients of $E_*$, which proves the claim. In particular $E_*(p)=0$ and
$\nabla E_*(p)=0$, since $L(p)=0$. In centered coordinates
$L(p+u)=\mathcal Bu$ with $\mathcal B=-aI+\mathcal U$, where $\mathcal U$ has
ones on the first superdiagonal. Hence $E_*(p+u)=u^{\mathsf T}H_*u$ with
$H_*=\mathcal B^{\mathsf T}\Theta\mathcal B$.

We have $\Theta=D_a\mathcal T D_a$, where
$D_a=\diag(a^s,a^{s-1},\ldots,1)$ has entries in $[1,\sqrt2)$ and
$\mathcal T$ has unit diagonal and entries $\frac12$ next to it. The
eigenvalues of $\mathcal T$ are $1+\cos(j\pi/(n+1))$, $1\le j\le n$, and
$1-\cos(\pi/(n+1))=2\sin^2(\pi/(2n+2))\ge2/(n+1)^2\ge1/(2n^2)$. Thus
$I/(2n^2)\preceq\Theta\preceq4I$. Next $\norm{\mathcal B}\le1+a<3$, and the
finite series $\mathcal B^{-1}=-a^{-1}\sum_{j=0}^{n-1}(\mathcal U/a)^j$ gives
$\norm{\mathcal B^{-1}}\le n$. Therefore
\begin{equation}
 \frac1{2n^4}I\preceq H_*\preceq36I.
 \label{eq:fields-prime-exposing-bounds}
\end{equation}
Each $P_j$ has quadratic part of norm at most one, and gradient at $p$ of
norm at most $5$: for example $\nabla P_1(p)$ has entries $-2$ and
$2a^{2s}<4$.

\paragraph{Chain residuals and their Jacobian.}
Put
\[
 C_0=x_0^2-x_s,\qquad C_1=x_0x_1-2,\qquad C_i=x_0x_i-x_1x_{i-1}\quad
 (2\le i\le s).
\]
They vanish at $p$, have no linear term in $x_0$, have quadratic parts of
norm at most one, and have gradients at $p$ of norm at most $5$. For a
direction $u$ put $h_i=u_i/p_i$, and divide the differential of each $C_i$ at
$p$ by the positive number $p_s$, $p_0p_1$, or $p_0p_i=p_1p_{i-1}$,
respectively, each in $[1,4]$. The result is the linear map
\[
 \varrho_0=2h_0-h_s,\qquad \varrho_1=h_0+h_1,\qquad
 \varrho_i=h_0+h_i-h_1-h_{i-1}\quad(2\le i\le s).
\]
With $\delta=h_1-h_0$ the last equations give
$h_i=h_0+i\delta+\sum_{j=2}^i\varrho_j$, and the first two then give
\[
 h_0=\frac{\varrho_0+s\varrho_1+\sum_{j=2}^s\varrho_j}{2s+1},\qquad
 \delta=\frac{\varrho_1-2\varrho_0-2\sum_{j=2}^s\varrho_j}{2s+1}.
\]
By Cauchy--Schwarz, $|h_0|\le\norm\varrho$ and
$|\delta|\le\norm\varrho\sqrt{4s+1}/(2s+1)\le\norm\varrho/\sqrt s$, so
$|h_i|\le(1+2\sqrt s)\norm\varrho$ and $\norm h\le3n\norm\varrho$. Since
$\norm u\le2\norm h$ and the row scaling divides by numbers at least one,
the Jacobian $J$ of $(C_0,\ldots,C_s)$ at $p$ satisfies
$\norm{Ju}\ge\norm u/(6n)$. Hence $J^{\mathsf T}J\succeq\nu^2I$ with
$\nu=1/(8n)$.

\paragraph{The integer quartic.}
Put
\[
 N_0=2^{20}n^8,\qquad t=N_0^{-1},\qquad \varepsilon=t^2,\qquad
 N_1=64(n+1)N_0^2,\qquad A=\sum_{j=0}^{s+1}\lfloor N_1w_j\rfloor P_j,
\]
and
\begin{equation}
 f_\ell=A^2+(N_1/N_0)^2\sum_{i=0}^sC_i^2-r^2,\qquad r=x_1x_s-2x_0 .
 \label{eq:fields-prime-quartic}
\end{equation}
The weight $-2$ is kept exactly, and $N_1/N_0=64(n+1)N_0$ is an integer, so
$f_\ell$ has integer coefficients. All of $A$, $C_i$, and $r$ vanish at $p$,
so $f_\ell(p)=0$ and $\nabla f_\ell(p)=0$.

Write $\bar A=A/N_1$. Each of the $n+1$ weights changes by at most $1/N_1$,
so $\bar A(p+u)=\ell_A^{\mathsf T}u+u^{\mathsf T}Hu$ with
\[
 \norm{\ell_A}\le\frac{5(n+1)}{N_1}<\frac\varepsilon8,\qquad
 \norm{H-H_*}\le\frac{n+1}{N_1}=\frac\varepsilon{64}.
\]
With \eqref{eq:fields-prime-exposing-bounds} this gives $H\succeq\mu I$ for
$\mu=1/(4n^4)$ and $\norm H\le40=:\Lambda$. Take $\beta=8$ and the $\nu$
above. The conditions \eqref{eq:algebraic-realization-eps} hold with $m=n$:
$\mu^2/(2n)=1/(32n^9)$, and since $\Lambda+n\beta=40+8n\le24n$,
\[
 \frac{\nu^2\mu^2}{36n(\Lambda+n\beta)^2}
 \ge\frac1{64n^2\cdot16n^8\cdot36n\cdot576n^2}=\frac1{21233664\,n^{13}},
\]
while $\varepsilon=2^{-40}n^{-16}$ is smaller than both and
$\norm{\ell_A}\le\varepsilon$. Lemma~\ref{lem:quartic-realization} applies to
$g=\bar A$ and the residuals $C_0,\ldots,C_s$. Here
$\theta\le3\sqrt n\,\varepsilon\cdot24n\cdot16n^8<1152n^{10}\varepsilon<1$,
the minimum in $\gamma$ equals $3/2$, and $\norm p\le2\sqrt n$ gives
$(1+\norm p)^2\le9n$. Hence the quartic
$f_{\rm lem}=(\bar A/(t\nu))^2+\sum_i(C_i/\nu)^2$ has a rational canonical Gram
$\Gamma_{f_{\rm lem}}\succeq\gamma I$ with $\gamma\ge1/(30n)$.

Since $f_\ell+r^2=N_1^2(\bar A^2+\varepsilon\sum_iC_i^2)
=N_1^2\varepsilon\nu^2f_{\rm lem}$, the canonical Gram of the representation
\eqref{eq:fields-prime-quartic} is
$\Gamma_{f_\ell}=N_1^2\varepsilon\nu^2\Gamma_{f_{\rm lem}}-\Gamma(r)$. At the
origin, $r$ has constant term zero, linear part $-2x_0$, and quadratic matrix
$T_r$ with $\norm{T_r}=\frac12$ and $\norm{T_r}_F^2=\frac12$, so the blocks of
$\Gamma(r)$ have norms at most $8$, $6\cdot2\cdot2^{-1/2}<9$, and $5$, and
$\norm{\Gamma(r)}<17$. As
$N_1^2\varepsilon\nu^2\gamma\ge64(n+1)^2N_0^2/(30n^3)\ge2N_0^2/n$,
\[
 \Gamma_{f_\ell}\succeq\bigl(2N_0^2/n-17\bigr)I\succeq I .
\]
This is an integer matrix, because the factors in
\eqref{eq:fields-prime-quartic} are integer quadratics with integer weights.
Since $\norm{w(x,v)}\ge\norm v$, it proves $\nabla^2f_\ell\succeq I$
everywhere, and with $f_\ell(p)=0$, $\nabla f_\ell(p)=0$, it gives
$f_\ell(x)\ge\frac12\norm{x-p}^2$. This proves
Theorem~\ref{thm:fields-prime}(b).

\paragraph{Size and construction.}
The quadratic $A$ has $O(n)$ monomials and coefficients bounded by $8N_1$,
and $N_1$ and $N_0$ are bounded by a polynomial in $n$. Hence $f_\ell$ has
$O(n^2)$ monomials with coefficients of magnitude polynomial in $n$, and the
canonical Gram entries are polynomial in these coefficients. This proves
Theorem~\ref{thm:fields-prime}(a). Each integer $\lfloor N_1w_j\rfloor$ for
$w_j=2^{e/\ell}$, $e\in\{2,4,\ldots,2s,\ell+1\}$, is the largest integer $y$
with $y^\ell\le N_1^\ell2^e$, found by binary search over $[0,4N_1]$ with
integers of $O(\ell\log n)$ bits. All other steps are arithmetic on $O(n^4)$
numbers of $O(\log n)$ bits, so the construction runs in time polynomial in
$\ell$. No number-field computation is needed.

\paragraph{The functional.}
For every quadratic $g=g_0+\sum_iv_ix_i+\sum_{i\le j}h_{ij}x_ix_j$ one
has $\varphi_\ell(g^2)=v_0^2+2g_0(v_s+h_{00})$, as in the proof of
Lemma~\ref{lem:fields-prime-functional}. The quadratics $A$ and $C_i$ are
rational and vanish at $p$, so $v_s+h_{00}=0$ by that lemma with $E=\Q$,
and they have $v_0=0$. Thus $\varphi_\ell(A^2)=\varphi_\ell(C_i^2)=0$, while
$r^2=x_1^2x_s^2-4x_0x_1x_s+4x_0^2$ gives $\varphi_\ell(r^2)=4$. Hence
$\varphi_\ell(f_\ell)=-4$, which is Theorem~\ref{thm:fields-prime}(c).

\begin{remark}[Other scales]
\label{rem:fields-prime-scales}
For every rational $\lambda\ge1$ the quartic
$f_\ell^{(\lambda)}=\lambda(f_\ell+r^2)-r^2$ has the canonical Gram
$\Gamma_{f_\ell}+(\lambda-1)N_1^2\varepsilon\nu^2\Gamma_{f_{\rm lem}}\succeq I$
and $\varphi_\ell(f_\ell^{(\lambda)})=-4$. The monomial count in
Theorem~\ref{thm:fields-prime}(a) and its conclusions (b)--(d) therefore hold
for $f_\ell^{(\lambda)}$, with the same proof. The quartic and this Hessian
Gram are rational, and are integer when $\lambda$ is integer. Given
$\lambda$, construction takes time polynomial in $\ell$ and its binary
length; coefficient and Gram-entry lengths are
$O(\log\ell+\langle\lambda\rangle)$.
\end{remark}

\subsection{The explicit ternary quartic}
\label{app:fields-ternary}

\paragraph{The convexity certificate.}
Let $c=(3/4,1,1/2)$ and $u=X-c$, and let
\[
 M_\star=\Gamma_c(A_0)+\sum_{i=0}^3\Gamma_c(r_i)-\Gamma_c(r_4),
\]
computed from \eqref{eq:algebraic-canonical} with the coordinates ordered
as $v_1,v_2,v_3$ followed by the products $u_kv_j$ in lexicographic order of
$(k,j)$. By
Lemma~\ref{lem:algebraic-gram}(a) it satisfies
$v^{\mathsf T}\nabla^2F_\star(c+u)v=(v,u\otimes v)^{\mathsf T}M_\star(v,u\otimes v)$.
The matrix $2(M_\star-I)$ is an integer matrix, and its leading principal
minors of orders $1,\ldots,12$ are
\begin{center}\small
\begin{tabular}{rl@{\qquad}rl}
1 & $98$ & 7 & $132268108059697184$\\
2 & $8714$ & 8 & $26367351626513255040$\\
3 & $1060526$ & 9 & $8176807399117693383936$\\
4 & $1508095300$ & 10 & $4716316942464353421986304$\\
5 & $505804465336$ & 11 & $2238789673795455159378824192$\\
6 & $388208815855184$ & 12 & $1231832682931718409670151748608$
\end{tabular}
\end{center}
All are positive, so $M_\star\succ I$ by Sylvester's criterion. Since
$\norm{(v,u\otimes v)}\ge\norm v$, this gives $\nabla^2F_\star\succeq I$ on
$\R^3$. By the covariance identity of Lemma~\ref{lem:algebraic-gram}(b),
$\Psi_c^{\mathsf T}M_\star\Psi_c$ is a rational positive definite full Hessian
Gram of $F_\star$ on $w(X,v)$. This finite
certificate is the only computational ingredient of
Example~\ref{ex:fields-ternary}; the remaining statements are proved below.

\paragraph{The zero and the vanishing basis.}
With $a^5=2$ and $p_\star=(a^4/2,a,a^2/2)$ one checks $xy=a^5/2=1$,
$x^2=yz=a^3/2$, $y^2=2z=a^2$, $2z^2=x=a^4/2$, and $2xz=y=a$ at $p_\star$, so
all $r_i$ vanish there. Let $E\subseteq\R$ be a field with $a\notin E$. The
monomials $1,y,z,yz,x$ take the values $1,a,a^2/2,a^3/2,a^4/2$, which are
linearly independent over $E$ by Lemma~\ref{lem:fields-tools}(d). The
evaluation map from the ten-dimensional space $E[X]_{\le2}$ to the
five-dimensional $E$-space $E(a)$ therefore has rank five, and
$\dim_E\mathcal I_2^E(p_\star)=5$. The monomials $xy$, $x^2$, $y^2$, $z^2$, $xz$
each occur in exactly one $r_i$, so the $r_i$ are independent and form a
basis.

\paragraph{The functional.}
Only $r_0$ has a constant term, namely $2$; only $r_2$ has a $z$ term,
namely $-4z$, and a $y^2$ term, namely $2y^2$; only $r_4$ has a $y$ term,
namely $-y$. Hence $[z](r_ir_j)\ne0$ only for $\{i,j\}=\{0,2\}$, where it is
$-8$, and $[y^2](r_ir_j)\ne0$ only for $\{i,j\}=\{0,2\}$, where it is $4$,
and for $i=j=4$, where it is $1$. With
$\varphi_\star=2[z]+4[y^2]$ this gives $\varphi_\star(r_0r_2)=-16+16=0$,
$\varphi_\star(r_4^2)=4$, and $\varphi_\star(r_ir_j)=0$ otherwise. For
$g=\sum_ic_ir_i$ we obtain $\varphi_\star(g^2)=4c_4^2\ge0$, and
$\varphi_\star(F_\star)=-4$ because $A_0$ has no $r_4$ component.
Lemma~\ref{lem:fields-tools}(c) excludes both certificate types over every
real field not containing $a$, and Lemma~\ref{lem:taylor-sos}(e) gives a
sum of squares over every field containing $a$.

\paragraph{Products.}
The leading forms of $r_0,\ldots,r_4$ are nonzero multiples of
$l_1=xy$, $l_2=x^2-yz$, $l_3=y^2$, $l_4=z^2$,
$l_5=xz$. Ten of the fifteen products $l_il_j$ are single monomials,
namely
\[
 x^2y^2,\ xy^3,\ xyz^2,\ x^2yz,\ y^4,\ y^2z^2,\ xy^2z,\ z^4,\ xz^3,\ x^2z^2 .
\]
The remaining five quartic monomials are
\[
 x^3y=l_1l_2+l_3l_5,\quad
 x^4=l_2^2+2l_1l_5-l_3l_4,\quad
 y^3z=l_1^2-l_2l_3,\quad
 yz^3=l_5^2-l_2l_4,\quad
 x^3z=l_2l_5+l_1l_4 .
\]
Thus the fifteen products span the fifteen-dimensional space of ternary
quartic forms and are independent. A linear relation among the products
$r_ir_j$ would give one among their quartic parts, so the $r_ir_j$ are
independent too, and the Gram matrix of $F_\star$ on $(r_0,\ldots,r_4)$ is
unique.

\subsection{Proof of Proposition~\ref{prop:fields-four}}
\label{app:fields-four}

\paragraph{The field.}
Let $K=\Q(\alpha)$, a real cubic field. If $\beta\in K$, write
$\beta=b_0+b_1\alpha+b_2\alpha^2$ with rational $b_i$. Then $\beta$ has
degree three and generates $K$, so its conjugates are the images of $\beta$
under the three embeddings of $K$, and $\operatorname{Tr}_{K/\Q}(\beta)=0$
and $\operatorname{Tr}_{K/\Q}(\beta^2)=0$, because the minimal polynomial
$T^3-5$ has zero coefficients of $T^2$ and $T$. Since
$\operatorname{Tr}(1)=3$ and $\operatorname{Tr}(\alpha)=\operatorname{Tr}(\alpha^2)=0$,
the first equation gives $b_0=0$. Then
$\beta^2=b_1^2\alpha^2+2b_1b_2\alpha^3+b_2^2\alpha^4
=4b_1b_2+b_1^2\alpha^2+2b_2^2\alpha$, and the second equation gives
$12b_1b_2=0$. If $b_2=0$ then $b_1^3=5/2$, and if $b_1=0$ then
$b_2^3=5/4$; neither is the cube of a rational. Hence $\beta\notin K$. The
polynomial $T^3-5$ has no root in the real field $K$, since its only real
root is $\beta$, so it is irreducible over $K$, and
$[\Q(\alpha,\beta):\Q]=9$, with basis $\alpha^i\beta^j$, $0\le i,j\le2$.

\paragraph{Vanishing quadratics and their products.}
The fifteen monomials of degree at most two in $(x,y,z,w)$ evaluate at
$p$ into $\Q(\alpha,\beta)$, and $1,x,y,z,w,xz,xw,yz,yw$ take the nine basis
values. The evaluation map has rank nine, so
$\dim_\Q\mathcal I_2^\Q(p)=6$. The six quadratics
\[
 g_1=x^2-y,\quad g_2=xy-2,\quad g_3=y^2-2x,\quad
 g_4=z^2-w,\quad g_5=zw-5,\quad g_6=w^2-5z
\]
vanish at $p$ and are independent, their pivot monomials
$x^2,xy,y^2,z^2,zw,w^2$ being distinct, so they form a basis. Products of two
of them either involve only one of the blocks $(x,y)$ and $(z,w)$ or have
degree at most two in each block. Neither kind contains the monomial
$yz^3$, so $[yz^3]P=0$ for every $P\in\mathcal W^\Q(p)$.

\paragraph{The perturbation.}
Let $h=y\,C(z,w)$ with $C=z^3+w^3/5-3zw+5$. At $(z,w)=(\beta,\beta^2)$ we have
$C=5+5-15+5=0$, $C_z=3z^2-3w=0$, and $C_w=\frac35w^2-3z=\frac35\cdot5\beta-3\beta=0$.
Hence $h(p)=0$ and $\nabla h(p)=0$, while $[yz^3]h=1$. The quartic $h$ has a
rational symmetric full Hessian Gram $N_h$: every monomial of its Hessian
biform has degree two in $v$ and at most two in $X$, and is therefore a
product of two entries of $w(X,v)$, so coefficients can be distributed over
the corresponding symmetric pairs of entries.

\paragraph{A strictly convex rational baseline.}
For $\rho>0$ let
$\mathcal E_\rho(X,Y)=\rho^2(X^2-Y)-\rho(XY-\rho^3)+(Y^2-\rho^3X)$. Expanding at
$X=\rho+u$, $Y=\rho^2+t$ gives $\mathcal E_\rho=\rho^2u^2-\rho ut+t^2$, a
positive definite form with determinant $3\rho^2/4$. Put
$E_*=\mathcal E_\alpha(x,y)+\mathcal E_\beta(z,w)
=\alpha^2g_1-\alpha g_2+g_3+\beta^2g_4-\beta g_5+g_6$. It vanishes at $p$ with
zero gradient and positive definite quadratic part. For rationals
$\tilde c_1,\ldots,\tilde c_4$ put
$E=\tilde c_1g_1-\tilde c_2g_2+g_3+\tilde c_3g_4-\tilde c_4g_5+g_6$; it
vanishes at $p$ exactly, and as $(\tilde c_i)\to(\alpha^2,\alpha,\beta^2,\beta)$
its quadratic part tends to that of $E_*$ and its gradient at $p$ tends to
zero. The residuals $g_1,g_2,g_4,g_5$ have Jacobian at $p$ block diagonal
with blocks $\left(\begin{smallmatrix}2\rho&-1\\\rho^2&\rho\end{smallmatrix}\right)$,
$\rho\in\{\alpha,\beta\}$, of determinant $3\rho^2\ne0$, and their quadratic
parts have norm at most one. Fix positive rationals $\mu$, $\Lambda$,
$\beta_0$, $\nu$ with $\mu$ at most half the least eigenvalue of the quadratic
part of $E_*$, $\Lambda$ at least twice its norm, $\beta_0$ at least the norms
of the residual gradients at $p$ (this is the bound called $\beta$ in
\eqref{eq:algebraic-realization-data}), and $\nu^2$ at most the least
eigenvalue of $J^{\mathsf T}J$ for their Jacobian $J$. Choose a rational $t$ such that
$\varepsilon=t^2$ satisfies the first part of
\eqref{eq:algebraic-realization-eps} with $m=n=4$, and then choose the
$\tilde c_i$ so close that the quadratic part of $E$ is at least $\mu I$ and
at most $\Lambda$ in norm, and $\norm{\nabla E(p)}\le\varepsilon$.
Lemma~\ref{lem:quartic-realization} applies to $g=E$ and the residuals
$g_1,g_2,g_4,g_5$: the rational sum of squares
$G=(E/(t\nu))^2+\sum_{j\in\{1,2,4,5\}}(g_j/\nu)^2$ has a rational positive
definite canonical Gram $M_G$. Moreover $G\in\mathcal W^\Q(p)$.

\paragraph{Conclusion.}
Let $\lambda$ be a rational number with
$\lambda\mu_{M_G}\ge1+\sum_{i,j}|(N_h)_{ij}|$, where $\mu_{M_G}$ is the
determinant--trace margin, and put $f=\lambda G+h$. The sum of absolute
entries bounds $\norm{N_h}$. Its full Hessian Gram
$\lambda M_G+N_h\succeq I$ is rational, so $\nabla^2f\succeq I$. Since
$f(p)=0$ and $\nabla f(p)=0$, the point $p$ is the unique minimizer and the
minimum is zero. The coefficient of $yz^3$ in $f$ is $\lambda\cdot0+1=1$, so
$f\notin\mathcal W^\Q(p)$. A sum of squares of rational polynomials would
consist of elements of $\mathcal I_2^\Q(p)$ by
Lemma~\ref{lem:fields-tools}(a,b) and would lie in $\mathcal W^\Q(p)$. Finally
$p\in\Q(\alpha,\beta)^4$, and Lemma~\ref{lem:taylor-sos}(e) gives a sum of
squares over that field.

\subsection{The tower baseline and the proof of Theorem~\ref{thm:fields-tower}}
\label{app:fields-tower}

\paragraph{One quintic gate.}
For $\alpha>0$ and $b=\alpha^5$ let $\bar r_1,\ldots,\bar r_5$ be the five
quadratics \eqref{eq:fields-tower-relations} of one gate with the constant
$b$ in place of $b_i$, and put
\[
 \mathcal E_\alpha=\bar r_5-\alpha\bar r_4+\alpha^4\bar r_1+\alpha^2\bar r_3 .
\]
Let $\ell_\alpha=(x-\alpha,\ y-\alpha x,\ z-\alpha y)$ and let $T_\alpha$ be the
symmetric tridiagonal matrix with diagonal $(\alpha^4,\alpha^2,1)$ and
off-diagonal entries $\alpha^3/2$ and $\alpha/2$. Expanding
$\ell_\alpha^{\mathsf T}T_\alpha\ell_\alpha$ gives
\[
 \alpha^4x^2+\alpha^2y^2-\alpha^2xz-\alpha yz+z^2-\alpha^5x-\alpha^4y+\alpha^6,
\]
because the coefficient of $xy$ is $\alpha^3-2\alpha^3+\alpha^3=0$; the same
expansion of $\mathcal E_\alpha$ with $b=\alpha^5$ gives the identical
polynomial. Hence
\begin{equation}
 \mathcal E_\alpha=\ell_\alpha^{\mathsf T}T_\alpha\ell_\alpha,\qquad
 \mathcal E_\alpha(\alpha+u_1,\alpha^2+u_2,\alpha^3+u_3)=u^{\mathsf T}H_\alpha u,
 \quad H_\alpha=\mathcal B_\alpha^{\mathsf T}T_\alpha\mathcal B_\alpha,
 \label{eq:fields-gate-exposing}
\end{equation}
where $\mathcal B_\alpha$ is lower bidiagonal with unit diagonal and
subdiagonal $-\alpha$. For $1\le\alpha\le2^{1/5}$ we have the following
bounds. $T_\alpha=D_\alpha\mathcal T_3D_\alpha$ with
$D_\alpha=\diag(\alpha^2,\alpha,1)\succeq I$, and the tridiagonal matrix
$\mathcal T_3$ has eigenvalues $1$ and $1\pm\sqrt2/2$, so
$T_\alpha\succeq(1-\sqrt2/2)I\succeq\frac14I$ and
$\norm{T_\alpha}\le(1+\sqrt2/2)\alpha^4\le3$. The inverse of
$\mathcal B_\alpha$ is lower triangular with entries $1,\alpha,\alpha^2$,
so $\norm{\mathcal B_\alpha^{-1}}^2\le3+2\alpha^2+\alpha^4<9$, while
$\norm{\mathcal B_\alpha}\le1+\alpha<2.2$. Hence
\begin{equation}
 \frac1{36}I\preceq H_\alpha\preceq16I\qquad(1\le\alpha\le2^{1/5}).
 \label{eq:fields-gate-bounds}
\end{equation}

\paragraph{Coupling the gates.}
For the tower put
$E_i=r_{i,5}-a_ir_{i,4}+a_i^4r_{i,1}+a_i^2r_{i,3}$, with $b_i=x_{i-1}$ for
$i\ge2$. Replacing the constant $a_i^5=a_{i-1}$ by $x_{i-1}$ in
$\bar r_5-a_i\bar r_4$ changes it by
$-(x_{i-1}-a_{i-1})x_i+a_i(x_{i-1}-a_{i-1})=-(x_{i-1}-a_{i-1})(x_i-a_i)$.
Writing $u=X-p_k$ and $u_i=(u_{i,1},u_{i,2},u_{i,3})$, we obtain from
\eqref{eq:fields-gate-exposing}
\[
 E_i(p_k+u)=u_i^{\mathsf T}H_{a_i}u_i-u_{i-1,1}u_{i,1}\quad(i\ge2),\qquad
 E_1(p_k+u)=u_1^{\mathsf T}H_{a_1}u_1 .
\]
Let $\rho=2^{-13}$ and $E_*=\sum_{i=1}^k\rho^{i-1}E_i$. By the
arithmetic-geometric mean inequality,
$\rho^{i-1}|u_{i-1,1}u_{i,1}|\le\frac{\rho^{i-1}}{144}u_{i,1}^2
+36\rho^{i-1}u_{i-1,1}^2$, and $36\rho^{i-1}\le\rho^{i-2}/144$ because
$\rho\le1/5184$. Summing and using \eqref{eq:fields-gate-bounds},
\[
 E_*(p_k+u)\ge\sum_i\rho^{i-1}\Bigl(\frac1{36}-\frac2{144}\Bigr)\norm{u_i}^2
 \ge\frac{\rho^{k-1}}{72}\norm u^2 .
\]
The quadratic part of $E_*$ also has norm at most $16+1=17$, the coupling
terms forming a matrix of norm at most one.

\paragraph{The rational exposing quadratic.}
Put $a_0=2$, $\vartheta=2^{-32k-15}$, and compute dyadic rationals
$\tilde a_1,\ldots,\tilde a_k\in[1,2]$ as follows: $\tilde a_0=2$, and
$\tilde a_i$ is obtained by bisection on $[1,2]$ for the increasing map
$\tau\mapsto\tau^5$, compared exactly with $\tilde a_{i-1}$, so that
$|\tilde a_i-\tilde a_{i-1}^{1/5}|\le\vartheta/2$. Since $\tau\mapsto\tau^{1/5}$ is
$\frac15$-Lipschitz on $[1,\infty)$, induction gives
$|\tilde a_i-a_i|\le\frac\vartheta2+\frac15|\tilde a_{i-1}-a_{i-1}|\le\vartheta$.
Each bisection uses $O(k)$ steps on rationals of $O(k)$ bits. Define
\[
 G_k=\sum_{i=1}^k\rho^{i-1}\bigl(r_{i,5}-\tilde a_ir_{i,4}
      +\tilde a_i^4r_{i,1}+\tilde a_i^2r_{i,3}\bigr).
\]
It is a rational quadratic, it vanishes at $p_k$ exactly, and
$[z_k^2]G_k=\rho^{k-1}>0$, since $z_k^2$ occurs only in $r_{k,5}$. The
difference $G_k-E_*$ is a combination of $r_{i,4},r_{i,1},r_{i,3}$ with
coefficients $\rho^{i-1}$ times $-(\tilde a_i-a_i)$,
$\tilde a_i^4-a_i^4$, and $\tilde a_i^2-a_i^2$, of absolute values at most
$\vartheta$, $8\vartheta$, and $3\vartheta$. These relations have quadratic parts of
norm at most one and gradients at $p_k$ of norm at most $2.3$, $2.6$, and
$3.3$. As $\sum_i\rho^{i-1}<2$, the quadratic part $H$ of $G_k$ and its
gradient $\ell_G$ at $p_k$ satisfy $\norm{H-H_{E_*}}\le24\vartheta$ and
$\norm{\ell_G}\le70\vartheta$.

\paragraph{The chain Jacobian.}
Let $J$ be the Jacobian at $p_k$ of the $3k$ chain residuals, rows ordered
by gates. The block of gate $i$ with respect to $(x_i,y_i,z_i)$ is
\[
 J_i=\begin{pmatrix}2\alpha&-1&0\\\alpha^2&\alpha&-1\\0&\alpha^3&\alpha^2\end{pmatrix},
 \quad
 J_i^{-1}=\frac1{5\alpha^4}\begin{pmatrix}
 2\alpha^3&\alpha^2&1\\-\alpha^4&2\alpha^3&2\alpha\\\alpha^5&-2\alpha^4&3\alpha^2
 \end{pmatrix},\quad\alpha=a_i,
\]
with $\det J_i=5\alpha^4$; for $\alpha\in[1,2^{1/5}]$ the Frobenius norm gives
$\norm{J_i^{-1}}<2$. The only other entries of $J$ are $-1$ in the row of
$r_{i,4}$ and the column of $x_{i-1}$. Writing $J=\mathcal D+\mathcal N$
with $\mathcal D=\diag(J_i)$, the matrix $\mathcal D^{-1}\mathcal N$ is
block strictly lower triangular with norm less than two, so
$J^{-1}=\sum_{j=0}^{k-1}(-\mathcal D^{-1}\mathcal N)^j\mathcal D^{-1}$ has
norm less than $2^{k+1}$. Hence $J^{\mathsf T}J\succeq\nu^2I$ with
$\nu=2^{-k-1}$. The chain residuals have quadratic parts of norm at most one
and gradients at $p_k$ of norm at most $3$.

\paragraph{The baseline.}
Put $\mu=\rho^{k-1}/144$, $\Lambda=18$, $\beta=3$, $n=3k$, $t=2^{-16k-4}$, and
$\varepsilon=t^2=2^{-32k-8}$. Then $24\vartheta\le\mu$, so $H\succeq\mu I$
and $\norm H\le\Lambda$, and $\norm{\ell_G}\le70\vartheta\le\varepsilon$. The
conditions \eqref{eq:algebraic-realization-eps} hold with $m=n$: since
$\Lambda+n\beta\le27k$, $20736\cdot78732<2^{31}$, and $k^3\le2^{3k}$,
\[
 \frac{\nu^2\mu^2}{36n(\Lambda+n\beta)^2}\ge
 \frac{2^{-2k-2}\,2^{-26(k-1)}}{20736\cdot78732\,k^3}
 \ge2^{-31k-7}\ge\varepsilon,
\]
and $\mu^2/(2n)$ is larger. By Lemma~\ref{lem:quartic-realization},
applied to $g=G_k$ and the $3k$ chain residuals,
\[
 F_{k,0}=\Bigl(\frac{G_k}{t\nu}\Bigr)^2
   +\sum_{i=1}^k\sum_{j\in\{1,2,4\}}\Bigl(\frac{r_{i,j}}\nu\Bigr)^2,
\]
an explicit sum of $3k+1$ squares of rational quadratics, has a rational
positive definite canonical Gram $M_k$, computed from these factors. In the
notation of Section~\ref{sec:fields-tower}, $\sigma_0=(t\nu)^{-2}$ and
$\sigma_1=\nu^{-2}$. Let
$\lambda_k=\lceil17/\mu_{M_k}\rceil$. All coefficients and the entries of
$M_k$ have $O(k)$ bits, the dimension of $M_k$ is $3k+9k^2$, and
$\mu_{M_k}$ and $\lambda_k$ are computed exactly in polynomial time with
polynomial bit length.

\paragraph{Proof of (a).}
The quadratic $r_{k,3}=X^{\mathsf T}TX$ is homogeneous, with $T$ supported on
the last gate, $\norm T=1$, and $\norm T_F^2=\frac32$. Its Gram
$\Gamma_0(r_{k,3})$ has only the quartic block, of norm at most
$8\cdot\frac32+4=16$. For rational $\lambda\ge\lambda_k$ the rational full
Hessian Gram $\lambda M_k-\Gamma_0(r_{k,3})$ of $f_{k,\lambda}$ is at least
$(\lambda\mu_{M_k}-16)I\succeq I$. All squares in $f_{k,\lambda}$ vanish at
$p_k$, so $f_{k,\lambda}(p_k)=0$ and $\nabla f_{k,\lambda}(p_k)=0$, and
strong convexity makes $p_k$ the unique minimizer.

\paragraph{Proof of (b).}
If $a_k\in E$, then $p_k\in E^{3k}$ since $a_i=a_k^{5^{k-i}}$, and
Lemma~\ref{lem:taylor-sos}(e) gives a sum of squares over $E$, hence also a
positive semidefinite Gram over $E$. Suppose $a_k\notin E$ and put
$L=E(a_k^5)$. By Lemma~\ref{lem:fields-tools}(d), $[L(a_k):L]=5$, so
$T^5-a_k^5$ is irreducible over $L$; for $k=1$, $L=E$. The coordinates
$a_i=(a_k^5)^{5^{k-1-i}}$, $i<k$, lie in $L$. For $P\in E[X]$ let
$P^\flat(x,y,z)=P(p_{<k},x,y,z)\in L[x,y,z]$, where $p_{<k}$ lists the
coordinates of the first $k-1$ gates. A sum of squares over $E$ restricts to a
sum of squares over $L$. If $f=z(X)^{\mathsf T}Qz(X)$ with $Q\succeq0$ over $E$,
then $z(p_{<k},V)=\Pi z'(V)$ for the monomial vector $z'$ in
$V=(x,y,z)$ and a matrix $\Pi$ over $L$, so
$f^\flat=z'^{\mathsf T}(\Pi^{\mathsf T}Q\Pi)z'$ with
$\Pi^{\mathsf T}Q\Pi\succeq0$ over $L$.

On the slice, every relation of an earlier gate becomes the constant zero,
and $r_{k,j}^\flat$ is the relation of one gate with $b=a_k^5$. Hence
\[
 f_{k,\lambda}^\flat=\lambda\Bigl(\sigma_0(G_k^\flat)^2+\sigma_1\sum_{j\in\{1,2,4\}}
 (r_{k,j}^\flat)^2\Bigr)-(r_{k,3}^\flat)^2 .
\]
Let $\bar p=(a_k,a_k^2,a_k^3)$. The monomials $1,x,y,z,xz$ take the values
$1,a_k,\ldots,a_k^4$, linearly independent over $L$, so
$\dim_L\mathcal I_2^L(\bar p)=5$; the relations $r^\flat_{k,j}$ vanish at $\bar p$ and have
the distinct pivot monomials $x^2,xy,y^2,yz,z^2$, so they form a basis. Their
leading forms $x^2,xy,y^2-xz,yz,z^2$ have products spanning all ternary
quartic forms: ten products are monomials, and
\[
 x^3z=\hat r_2^2-\hat r_1\hat r_3,\ \ xy^3=\hat r_2\hat r_3+\hat r_1\hat r_4,\ \
 y^3z=\hat r_3\hat r_4+\hat r_2\hat r_5,\ \ xz^3=\hat r_4^2-\hat r_3\hat r_5,\ \
 y^4=\hat r_3^2+2\hat r_2\hat r_4-\hat r_1\hat r_5,
\]
where $\hat r_j$ is the leading form of $r^\flat_{k,j}$. So the fifteen products
$r^\flat_{k,j}r^\flat_{k,j'}$ are linearly independent over $\R$, and the map
$S\mapsto(r^\flat)^{\mathsf T}Sr^\flat$ on real symmetric $5\times5$ matrices is
injective. Write $G_k^\flat=\sum_jc_jr^\flat_{k,j}$ with $c_j\in L$; then
$c_5=[z^2]G_k^\flat=[z_k^2]G_k>0$. The slice is represented by
\[
 S_0=\lambda\bigl(\sigma_0cc^{\mathsf T}+\sigma_1(e_1e_1^{\mathsf T}
   +e_2e_2^{\mathsf T}+e_4e_4^{\mathsf T})\bigr)-e_3e_3^{\mathsf T},
 \qquad
 y^{\mathsf T}S_0y=-c_5^2<0\ \text{ for }y=(0,0,c_5,0,-c_3).
\]
A sum of squares of $f_{k,\lambda}$ over $E$ restricts to
$f^\flat=\sum_j(d_j^{\mathsf T}r^\flat)^2$ with $d_j\in L^5$, by
Lemma~\ref{lem:fields-tools}(a,b), so injectivity gives
$S_0=\sum_jd_jd_j^{\mathsf T}\succeq0$, a contradiction. A positive
semidefinite Gram over $E$ restricts to one over $L$, and
Lemma~\ref{lem:fields-tools}(b) gives $f^\flat=(r^\flat)^{\mathsf T}S'r^\flat$ with
$S'\succeq0$, so $S'=S_0$, again a contradiction. Finally $T^{5^k}-2$ is
irreducible over $\Q$ by Eisenstein's criterion at two, so $\Q(a_k)$ has
degree $5^k$.

\subsection{Proof of Theorem~\ref{thm:fields-individual}}
\label{app:fields-individual}

Let $D=\ell^k$, $a=2^{1/D}$, $\zeta=e^{2\pi\mathrm i/D}$, $Z=\Q(\zeta)$, and
$Z^+=\Q(\zeta+\zeta^{-1})\subseteq\R$. The field $Z^+$ is a subfield of the
abelian extension $Z/\Q$, hence Galois over $\Q$, and it is totally real:
every embedding sends $\zeta+\zeta^{-1}$ to $2\cos(2\pi j/D)$.

\emph{Step 1: $[Z^+(a):Z^+]=D$.} By Lemma~\ref{lem:fields-tools}(d) the
minimal polynomial of $a$ over $Z^+$ is $T^r-a^r$ with $r\mid D$ and
$a^r\in Z^+$. If $r<D$, then $a^r=2^{1/(D/r)}$ has minimal polynomial
$T^{D/r}-2$ over $\Q$, by Eisenstein's criterion, and this polynomial has
nonreal roots because $D/r\ge3$ is odd. As $Z^+$ is normal over $\Q$, it
would contain these nonreal conjugates, contradicting $Z^+\subseteq\R$.

\emph{Step 2: $[Z(a):Z]=D$.} The real field $Z^+(a)$ does not contain
$\zeta$, which is a root of $T^2-(\zeta+\zeta^{-1})T+1\in Z^+[T]$. Hence
$[Z^+(a,\zeta):Z^+(a)]=2$ and likewise $[Z:Z^+]=2$. Since
$Z(a)=Z^+(a,\zeta)$, comparing $[Z(a):Z^+]=2D$ computed through $Z^+(a)$
and through $Z$ gives the claim.

\emph{Step 3: an automorphism of order $D$.} The field $S=Z(a)$ is the
splitting field of $T^D-2$ over $\Q$. By Step 2, $T^D-2$ is irreducible over
$Z$, so some $\sigma\in\operatorname{Gal}(S/Z)$ maps $a$ to $\zeta a$. Since
$\sigma$ fixes $\zeta$, $\sigma^j(a)=\zeta^ja$, and $\sigma$ has order
exactly $D$.

\emph{Step 4: lifting.} Let $K'$ be the compositum of the normal closures of
the fields $\Q(\beta_i)$; it is finite and Galois over $\Q$. It contains
$a$, hence all conjugates of $a$, and therefore $S$. Restriction
$\operatorname{Gal}(K'/\Q)\to\operatorname{Gal}(S/\Q)$ is surjective, so
some $\tau\in\operatorname{Gal}(K'/\Q)$ restricts to $\sigma$, and $D$
divides the order of $\tau$. On the other hand $\operatorname{Gal}(K'/\Q)$
acts faithfully on the union of the sets of conjugates of the $\beta_i$,
since these generate $K'$, and preserves each set. This embeds the group
into $\prod_i\mathfrak S_{d_i}$, $d_i=[\Q(\beta_i):\Q]$. If every $d_i<D$,
every cycle of every component permutation has length less than $\ell^k$,
so the $\ell$-adic valuation of the order of every element of the product is
at most $k-1$. This contradicts $\ell^k\mid\operatorname{ord}(\tau)$.

\emph{The certificate consequence.} If the coefficients of the square
factors of a sum of squares of $f_{k,\lambda}$ are algebraic, the real field
$E$ they generate admits a sum of squares, so $a_k\in E$ by
Theorem~\ref{thm:fields-tower}, and the lemma applies to the list of these
coefficients with $\ell=5$. The same argument applies to the entries of a
positive semidefinite Gram matrix with algebraic entries. A polynomial with a
root of degree $5^k$ has degree at least $5^k$, so a dense list of any
annihilating polynomial of that coefficient has at least $5^k+1$ entries.

\subsection{Proof of Theorem~\ref{thm:fields-compact}}
\label{app:fields-compact}

Let $M'=\lambda M_k-\Gamma_0(r_{k,3})\succeq I$ be the rational full Hessian
Gram of $f=f_{k,\lambda}$ from the proof of Theorem~\ref{thm:fields-tower}(a).

(a) Lemma~\ref{lem:taylor-sos}(e) writes $f$ as a sum of squares of
quadratics whose coefficients are rational polynomials of degree at most two
in the coordinates of $p_k$. Their number and the bit lengths of these
rational polynomials are polynomial in $L_\lambda$, since the entries of $M'$
have length polynomial in $L_\lambda$ and rational elimination of $M'$ has
polynomial bit complexity, and a positive rational $u/v$ is a sum of at most
$2\log_2(2uv)$ rational squares, by the binary expansion of $uv$ and the
identity $u/v=uv/v^2$ (Lemma~\ref{lem:models-rational-squares}). The circuit
computes $a_i$ by fifth roots in sequence, $a_i^2,a_i^3$ by two
multiplications per gate, and every coefficient by evaluating its rational
polynomial.

(b) Let $Y$ be a second vector of $3k$ variables, $u=X-Y$,
$\mathcal A=(u,Y\otimes u)$, $\mathcal B=(0,u\otimes u)$, and
\[
 s(X,Y)=\tfrac12\bigl(\mathcal A+\tfrac13\mathcal B\bigr)^{\mathsf T}M'
   \bigl(\mathcal A+\tfrac13\mathcal B\bigr)+\tfrac1{36}\mathcal B^{\mathsf T}M'\mathcal B .
\]
Writing $M'=\sum_jc_jc_j^{\mathsf T}$ with rational $c_j$ makes $s$ an explicit
rational sum of squares of degree four. Since
$w(Y+tu,u)=\mathcal A+t\mathcal B$, the integral identity of
Lemma~\ref{lem:taylor-sos}(a) and Taylor's formula with integral remainder,
both identities between polynomials in $(X,Y)$, give
\[
 f(X)=f(Y)+\nabla f(Y)^{\mathsf T}(X-Y)+s(X,Y).
\]
With $f=\sum_j\sigma_jg_j^2$ and
$H_j=\sigma_j\bigl(g_j(Y)+2\nabla g_j(Y)^{\mathsf T}(X-Y)\bigr)$ we have
$f(Y)+\nabla f(Y)^{\mathsf T}(X-Y)=\sum_jH_jg_j(Y)$, which is the identity.
Each $H_j$ has degree at most two. At $Y=p_k$ all $g_j(Y)$ vanish, so
$f(X)=s(X,p_k)\ge0$ for every real $X$. For the second identity, the exact
relations
\[
 r_{i,3}=-y_ir_{i,1}+x_ir_{i,2},\qquad r_{i,5}=-z_ir_{i,2}+x_ir_{i,4}
\]
express $r_{k,3}$ and $G_k$ as combinations of chain residuals with affine
rational multipliers. Substituting them multiplies $H_j$ by affine
polynomials, which gives multipliers of degree at most three and an identity
of degree at most five.

\subsection{Proof of Theorem~\ref{thm:fields-tower-spaces}}
\label{app:fields-spaces}

\paragraph{(a) The quadratic relations.}
Each relation $r_{i,j}$ has a pivot monomial, namely
$x_i^2,x_iy_i,y_i^2,y_iz_i,z_i^2$, and its other monomials lie outside the
set of pivots: they are $y_i$, $z_i$, $x_iz_i$, the constant or $x_{i-1}$, and
$x_1$ or $x_{i-1}x_i$. The monomials of degree at most two that are not
pivots are the constant, the four monomials $x_i,y_i,z_i,x_iz_i$ of each gate,
and the nine products of a variable of gate $i$ with a variable of gate $j$,
$i<j$. Let $a=a_k$, so that $a_i=a^{5^{k-i}}$. These monomials evaluate at
$p_k$ to $a^e$ with $0\le e<5^k$, where the base-five expansion of $e$ is,
respectively, zero; a single digit in $\{1,2,3,4\}$ at position $k-i$; or
two digits in $\{1,2,3\}$ at positions $k-i$ and $k-j$. Distinct monomials
give distinct expansions. By Eisenstein's criterion $T^{5^k}-2$ is
irreducible, so these values are linearly independent over $\Q$. Given
$g\in\mathcal I_2^\Q(p_k)$, subtract multiples of the relations to remove all
pivots. The remainder is a combination of non-pivot monomials vanishing at
$p_k$, hence zero. So the relations span, and their distinct pivots make them
independent. Since $1,x_i,y_i,z_i$ are non-pivot monomials, no nonzero
rational affine polynomial vanishes at $p_k$. It follows that
$\mathcal J_2^\Q(p_k)=\{0\}$: the partial derivatives of an element are
rational affine polynomials vanishing at $p_k$, so the element is a constant
equal to its value zero.

\paragraph{(b) Products.}
Induct on $k$. Let $V=(x_k,y_k,z_k)$ and grade polynomials by their degree
in $V$. Relations of earlier gates have $V$-degree zero, and
$r_{k,j}=\hat r_j+(\text{terms of }V\text{-degree at most one})$ with the
leading forms $\hat r=(x_k^2,x_ky_k,y_k^2-x_kz_k,y_kz_k,z_k^2)$. In a real
linear relation among the products, the part of $V$-degree four involves
only products of two last-gate relations and equals the same combination of
the products $\hat r_j\hat r_{j'}$; these are independent (proof of
Theorem~\ref{thm:fields-tower}(b)), so those coefficients vanish. The part of
$V$-degree two then reads $\sum_j\hat r_jP_j=0$, where $P_j$ is the
combination of earlier relations multiplying $r_{k,j}$. The forms $\hat r_j$
have linearly independent coefficient vectors, so they remain independent
over the field of rational functions in the earlier variables, and every
$P_j=0$. By (a) for the earlier gates, the earlier relations are independent,
so the coefficients of all mixed products vanish. The remaining relation
involves earlier gates only and is trivial by induction; for $k=1$ only the
local products occur.

\paragraph{(c) A local interpolation lemma.}
Let $L\subseteq\R$ be a field, $b\in L\setminus\{0\}$, $T^5-b$ irreducible
over $L$, and $\alpha$ its real root. We claim that the $L$-linear map
\[
 \mathcal H:L[x,y,z]_{\le3}\to L(\alpha)^4,\qquad
 P\mapsto(P,P_x,P_y,P_z)(\alpha,\alpha^2,\alpha^3)
\]
is bijective. Both sides have dimension twenty over $L$. Compose with the
invertible map multiplying the last three components by
$\alpha,\alpha^2,\alpha^3$. The composite sends $x^iy^jz^h$ to
$\alpha^w(1,i,j,h)$, where $w=i+2j+3h=5t+\varrho$ with $0\le\varrho\le4$,
and $\alpha^w=b^t\alpha^\varrho$. In the $L$-basis $\alpha^\varrho e_c$ of
$L(\alpha)^4$, the composite is block diagonal over the residue $\varrho$.
Each residue class contains four monomials:
\begin{center}\small
\begin{tabular}{cll}
$\varrho$ & monomials & determinant of the columns $(1,i,j,h)$\\\hline
0 & $1,\ yz,\ x^2z,\ xy^2$ & $5$\\
1 & $x,\ z^2,\ xyz,\ y^3$ & $5$\\
2 & $y,\ x^2,\ xz^2,\ y^2z$ & $-5$\\
3 & $z,\ xy,\ x^3,\ yz^2$ & $-5$\\
4 & $xz,\ y^2,\ x^2y,\ z^3$ & $-5$
\end{tabular}
\end{center}
For instance, for $\varrho=0$ the columns are $(1,0,0,0)$, $(1,0,1,1)$,
$(1,2,0,1)$, $(1,1,2,0)$, and the determinant equals that of
$\left(\begin{smallmatrix}0&2&1\\1&0&2\\1&1&0\end{smallmatrix}\right)$, which
is $5$. The block determinants are these integers times nonzero powers of
$b$, so the map is bijective.

\paragraph{(c) Stationary cubics.}
We show $\mathcal J_3^\Q(p_k)=\{0\}$ by induction. For $k=1$ apply the local
lemma with $L=\Q$ and $b=2$. Let $k\ge2$, $L=\Q(a_{k-1})$, $b=a_{k-1}$, and
$\alpha=a_k$; then $[L(\alpha):L]=5$, so $T^5-b$ is irreducible over $L$.
Write $U$ for the earlier variables and $F=\sum_\mu c_\mu(U)\mu(V)$ for
$F\in\mathcal J_3^\Q(p_k)$, summing over monomials $\mu$ in $V$. The
polynomial $F(p_{k-1},V)$ has degree at most three in $V$, coefficients in
$L$, and zero value and $V$-gradient at $(a_k,a_k^2,a_k^3)$, so it is zero.
Hence every $c_\mu$ vanishes at $p_{k-1}$. For $\deg\mu\ge2$ the coefficient
$c_\mu$ is affine, hence zero by (a). Thus
$F=h_0+x_kh_1+y_kh_2+z_kh_3$ with $\deg h_0\le3$, $\deg h_j\le2$, all
vanishing at $p_{k-1}$. For an earlier variable $\xi$,
$0=\partial_\xi F(p_k)=\partial_\xi h_0+a_k\partial_\xi h_1+a_k^2\partial_\xi h_2
+a_k^3\partial_\xi h_3$ evaluated at $p_{k-1}$, with coefficients in $L$;
independence of $1,a_k,a_k^2,a_k^3$ over $L$ makes each
$\partial_\xi h_j(p_{k-1})=0$. Thus
$h_1,h_2,h_3\in\mathcal J_2^\Q(p_{k-1})=\{0\}$ and
$h_0\in\mathcal J_3^\Q(p_{k-1})=\{0\}$.

\paragraph{(c) Stationary quartics.}
The inclusion $\mathcal W^\Q(p_k)\subseteq\mathcal J_4^\Q(p_k)$ follows from
the product rule; uniqueness of $S$ is (b). For the reverse inclusion let
$F\in\mathcal J_4^\Q(p_k)$. For $k=1$, subtract a rational combination of the
products $r_{1,j}r_{1,j'}$ with the same quartic part, which exists because
the leading products span all ternary quartic forms; the remainder lies in
$\mathcal J_3^\Q(p_1)=\{0\}$. For $k\ge2$ proceed in four steps.

\emph{Step 1.} The part of $F$ of $V$-degree four has rational constant
coefficients. Subtract the rational combination of products of last-gate
relations with this $V$-degree-four part. The remainder $F'$ lies in
$\mathcal J_4^\Q(p_k)$ and has $V$-degree at most three.

\emph{Step 2.} As in the cubic case, $F'(p_{k-1},V)=0$, so all coefficients
$c_\mu$ of $F'$ vanish at $p_{k-1}$. Coefficients of $V$-monomials of degree
three are affine and hence zero. Coefficients of $V$-monomials of degree two
have degree at most two and lie in $\mathcal I_2^\Q(p_{k-1})$.

\emph{Step 3.} We claim $c_{y^2}+c_{xz}=0$, where these are the coefficients
of $y_k^2$ and $x_kz_k$. For an earlier variable $\xi$, the polynomial
$\Phi_\xi(V)=\sum_\mu\partial_\xi c_\mu(p_{k-1})\mu(V)$ has degree at most two,
coefficients in $L$, and value $\partial_\xi F'(p_k)=0$ at
$(a_k,a_k^2,a_k^3)$. By the slice computation in the proof of
Theorem~\ref{thm:fields-tower}(b), it is an $L$-combination of the five
relations of one gate, and each of them has the same coefficient on $y^2$ as
minus its coefficient on $xz$. Hence $\partial_\xi(c_{y^2}+c_{xz})(p_{k-1})=0$
for every earlier variable $\xi$. Also $(c_{y^2}+c_{xz})(p_{k-1})=0$, so
$c_{y^2}+c_{xz}\in\mathcal J_2^\Q(p_{k-1})=\{0\}$.

\emph{Step 4.} By Step 3 the $V$-degree-two part of $F'$ equals
$\sum_j\hat r_jP_j$ with $P_1=c_{x^2}$, $P_2=c_{xy}$, $P_3=c_{y^2}$,
$P_4=c_{yz}$, $P_5=c_{z^2}$, all in $\mathcal I_2^\Q(p_{k-1})$, which by (a)
is spanned by the earlier relations. Subtracting $\sum_jr_{k,j}P_j$, a
combination of products of relations, leaves
$F''=h_0+x_kh_1+y_kh_2+z_kh_3\in\mathcal J_4^\Q(p_k)$ with $\deg h_0\le4$
and $\deg h_j\le3$. The terms $-x_{k-1}P_4$ and $-x_{k-1}x_kP_5$ of the
subtracted products only contribute to these coefficients. The value and
$V$-gradient of $F''$ at $p_k$ give $h_j(p_{k-1})=0$ for all $j$, and the
earlier derivatives give $\nabla h_j(p_{k-1})=0$ as in the cubic case. Hence
$h_1,h_2,h_3\in\mathcal J_3^\Q(p_{k-1})=\{0\}$, and
$h_0\in\mathcal J_4^\Q(p_{k-1})$ is a combination of products of earlier
relations by induction. This proves $F\in\mathcal W^\Q(p_k)$.

\subsection{Proof of Proposition~\ref{prop:fields-cyclic}}
\label{app:fields-cyclic}

Put $B=-2$ and $a=2^{1/d_n}$, so $e_i=(B^i-1)/3$, $e_0=0$, $p_i=a^{e_i}$,
and $3d_n=|B^m-1|$. Since
$2e_i-e_{i+1}-e_{i+2}=B^i(2-B-B^2)/3=0$ for $0\le i\le m-3$, and the two
wrapped indices give $\pm d_n$, each cyclic binomial vanishes at $p$; this is
the vanishing identity of the cyclic construction. Homogenize with $X_0=1$
and index monomials of degree at most two as $X_iX_j$,
$0\le i\le j\le m-1$. Such a monomial evaluates to
$a^{e_i+e_j}=2^{\lfloor(e_i+e_j)/d_n\rfloor}a^{(e_i+e_j)\bmod d_n}$, and
$1,a,\ldots,a^{d_n-1}$ are linearly independent over $\Q$. Group the monomials
into classes with equal residue of $e_i+e_j$ modulo $d_n$. The evaluation map
has rank equal to the number of classes, so $\dim\mathcal I_2^\Q(p)$ equals
the number of monomials minus the number of classes.

Two monomials $X_iX_j$ and $X_uX_v$ lie in one class exactly when the integer
$B^i+B^j-B^u-B^v=3(e_i+e_j-e_u-e_v)$ is divisible by $3d_n$, that is, when
$B^i+B^j\equiv B^u+B^v\pmod{B^m-1}$. Suppose the two monomials are distinct.
At most four indices occur, and $m\ge5$, so some index is unused.
Multiplication by a power of $B$, a unit modulo $B^m-1$ with $B^m\equiv1$,
shifts all indices cyclically, and we may assume that the unused index is
$m-1$. Then all four exponents are at most $m-2$. A term has modulus $2^{m-2}$
only if its exponent is $m-2$, and then it equals $B^{m-2}$; the other terms
have modulus at most $2^{m-3}$. If $n_+$ of the two terms $B^i,B^j$ and
$n_-$ of the two terms $B^u,B^v$ have exponent $m-2$, then
\[
 |B^i+B^j-B^u-B^v|\le|n_+-n_-|\,2^{m-2}+(4-n_+-n_-)\,2^{m-3}
 \le3\cdot2^{m-2}<2^m-1\le|B^m-1| .
\]
So the congruence is an equality of integers. If $i\ne j$ and $u\ne v$, the $2$-adic valuations give
$\min\{i,j\}=\min\{u,v\}$, and cancelling that term gives
$\max\{i,j\}=\max\{u,v\}$, so the monomials coincide. If $i=j$ and $u=v$,
then $i=u$. In the remaining case $2B^\varrho=B^u+B^v$ with $u<v$, the
valuations give $u=\varrho+1$, and then $B^v=B^\varrho(2-B)=4B^\varrho=B^{\varrho+2}$.
Shifting back, the only nontrivial classes are the pairs
$\{X_\varrho^2,X_{\varrho+1}X_{\varrho+2}\}$, indices modulo $m$. These $m$
pairs are distinct and pairwise disjoint, and a square has a unique partner,
so every class has at most two elements and exactly $m$ classes have two.
Hence $\dim\mathcal I_2^\Q(p)=m=n+1$, and the $m$ cyclic binomials, which
vanish at $p$ and involve distinct squares, form a basis. The last assertion
of the proposition is
Corollary~\ref{cor:fields-product-independence}, applied with the globally
convex rational sum of squares of Theorem~\ref{thm:algebraic-cyclic}, which
has minimum zero at $p$ and a positive definite Hessian there, and has degree
exactly four because its positive definite canonical Hessian Gram makes its
quartic form positive definite.
